\documentclass[landscape,12pt]{article}
\usepackage[letterpaper,margin=0.6in]{geometry}
\usepackage[T1]{fontenc}
\usepackage{lmodern}
\usepackage{microtype}
\usepackage{amsmath,amssymb}
\usepackage{booktabs}
\usepackage{array}
\usepackage{xcolor}
\usepackage{titlesec}
\usepackage{enumitem}
\setlist[itemize]{leftmargin=1.3em,itemsep=2pt,topsep=3pt}
\pagestyle{empty}
\emergencystretch=3em

\definecolor{model}{HTML}{0D47A1}
\definecolor{bench}{HTML}{1B5E20}
\definecolor{phys}{HTML}{B71C1C}
\definecolor{prior}{HTML}{8D6E00}

\newcommand{\pagehead}[2]{%
  {\Large\bfseries #1}\\[1pt]
  {\small\itshape #2}\\[-4pt]
  \rule{\linewidth}{0.6pt}\\[4pt]}

\begin{document}

% ================================================== PAGE 1
\pagehead{What the model computes, and what it does not}%
{Hunter Kinder --- accompanying the preprint Caixia Wan forwarded. One page on the model, one on
why I am writing to you, one of questions I cannot answer myself.}

\vspace{2mm}
\begin{minipage}[t]{0.485\linewidth}
{\large\bfseries\color{model} The object}\\[3pt]
A biofilm is represented as biomass parcels on a 3-D lattice, with a photon-transport calculation
(OpenMC) supplying a static dose-rate field and a radial PDE carrying a mobile contaminant into
the film. The transport runs \emph{once}; it does not read the biology back. That is a declared
design choice, not a physical law, and the manuscript says where it weakens.

\vspace{2mm}
{\large\bfseries\color{model} The question it is built to ask}\\[3pt]
Given a hydrated layer of known composition on a surface in a known field: what dose does the
layer receive, how much does it hold, and does its presence perturb the field around it?

\vspace{2mm}
{\large\bfseries\color{phys} What it is not}\\[3pt]
\begin{itemize}
  \item \textbf{Not calibrated.} Every parameter in the table is a literature-anchored prior, and
        a literature prior is not a calibration. The manuscript declares seven of twenty-six
        parameter rows inadmissible on the basis this model can accept.
  \item \textbf{Not a radiotrophy result.} Radiotrophy is not established for any of the seven
        species modelled, and the paper says so in its own words rather than importing the label.
        One melanin coefficient is retained with its sign declared wrong.
  \item \textbf{Not a measurement of anything.} The one radial result reported runs opposite to
        its own coefficient and sits within 1.4 standard errors of the seeding null. It is
        reported that way.
\end{itemize}
\end{minipage}\hfill
\begin{minipage}[t]{0.485\linewidth}
{\large\bfseries\color{prior} The gap that matters for you}\\[3pt]
Dose and dose rate are not interchangeable here, and the gap is not marginal. A reactor
irradiation delivers of order Gy\,min$^{-1}$. The environmental settings the paper cites as
motivation --- Chernobyl, Hanford --- are closer to mGy\,yr$^{-1}$. That is roughly ten orders of
magnitude, and a sensitivity constant calibrated in one regime does not transfer to the other by
construction.

\vspace{2mm}
That is the model's worst structural problem, and it is why I am not writing to you about
environmental remediation.

\vspace{4mm}
\noindent\rule{\linewidth}{0.4pt}\\[3pt]
{\small\textbf{One correction, since you may have version 1.0.} Figures 1 and 2 in that version
carry text inside the images that the surrounding prose retracts --- Figure 2 annotates two
species as ``radiotrophic (melanin-mediated energy gain)'' beneath a caption disowning exactly
that. The figure generator had been corrected two weeks earlier; the committed images had not
been regenerated, and nothing in the test suite could open a figure to notice. Version 1.1
replaces both. No number changed. I mention it because it is the paper's own argument happening
to the paper, and you would have found it on page 8.}
\end{minipage}

\newpage
% ================================================== PAGE 2
\pagehead{Why I am writing to a materials scientist}%
{The overlap is not bioremediation. It is what grows on the alloys, in the water, and what dose it
takes.}

\vspace{2mm}
\begin{minipage}[t]{0.485\linewidth}
{\large\bfseries\color{bench} The empirical fact}\\[3pt]
Sarr\'{o} and colleagues submerged a bioreactor carrying stainless steel and titanium coupons in
a spent nuclear fuel pool at a Spanish nuclear power plant. Biofilm developed on both materials.
They measured its radioactivity by gamma-ray spectrometry and found it had retained radionuclides
from the pool water, $^{60}$Co in particular.\footnote{\textit{J.\ Ind.\ Microbiol.\ Biotechnol.}
34(6):433--441, 2007. doi:10.1007/s10295-007-0215-7}

\vspace{2mm}
Karley and colleagues isolated six bacterial species from spent-fuel-pool water with $D_{10}$
values from 248~Gy to 2~kGy, all biofilm-formers, and measured Co and Ni removal up to
3.8~$\mu$g per mg of biomass.\footnote{\textit{Environ.\ Sci.\ Pollut.\ Res.}
25(21):20518--20526, 2018. doi:10.1007/s11356-017-0376-5} A later study of the same isolates
found \emph{dead} biomass removed both ions as well --- which is what a sorption mechanism
predicts and a metabolic one does not.\footnote{\textit{Environ.\ Monit.\ Assess.} 195(6):731,
2023. doi:10.1007/s10661-023-11266-x} One of those isolates is \textit{B.\ subtilis}, which is
already one of the seven species in the model.

\vspace{2mm}
{\large\bfseries\color{phys} What nobody in that literature computes}\\[3pt]
What dose the biofilm itself receives, and whether a hydrated layer 50--200~$\mu$m thick
perturbs the local field at the metal surface. The microbiology papers measure what the biofilm
holds. The transport question next to it is not asked.
\end{minipage}\hfill
\begin{minipage}[t]{0.485\linewidth}
{\large\bfseries\color{model} Why this regime and not the other}\\[3pt]
A facility is the setting where the model's inputs are obtainable rather than assumed:

\begin{itemize}
  \item \textbf{Bounded geometry} --- a coupon, not a heterogeneous subsurface.
  \item \textbf{Characterised source} --- pool water and fuel inventory are known, not inferred.
  \item \textbf{A measurable metal inventory in the film} --- the papers above already do this.
  \item \textbf{And a dose rate in the regime the $D_{10}$ priors were actually measured in},
        which closes the ten-order-of-magnitude gap on the previous page rather than talking
        around it.
\end{itemize}

\vspace{3mm}
{\large\bfseries\color{prior} Being straight about what this is}\\[3pt]
This is a reframing of application, not of content. The model does not change; what changes is
who the answer is for. I am not going to tell you that this work was aimed at your field, because
it was not --- it was aimed at contaminated sites, and the regime argument above is what moved
it. You work on cladding alloys in reactor water chemistry, and pitting is where
microbiologically influenced corrosion lives. The overlap is a question about what sits on the
metal, not a claim about the metal.
\end{minipage}

\newpage
% ================================================== PAGE 3
\pagehead{What the measurement needs}%
{Questions, not an inventory. You know what is available at MURR and I do not.}

\vspace{2mm}
\begin{minipage}[t]{0.45\linewidth}
{\large\bfseries\color{bench} The one thing the model cannot proceed without}\\[3pt]
A \textbf{closed composition for a hydrated biofilm}, mass fractions summing to unity, because
that is what a photon-transport code consumes. It needs three separable pieces:

\begin{itemize}
  \item \textbf{Water fraction} --- hydrated versus dry mass, gravimetric, with the drying
        protocol fixed and its own variance measured first.
  \item \textbf{The light elements} --- C, H, N, O, which are most of the mass and which I
        believe activation methods cannot see.
  \item \textbf{Trace metals, P, and any actinides} --- the part that binds, and the part the
        pool-water literature already measures.
\end{itemize}

\vspace{2mm}
The third is the piece I would most expect a reactor facility to be good at, and the first two
are the ones I would expect to go elsewhere --- but that is my reading of the technique, not of
your facility.
\end{minipage}\hfill
\begin{minipage}[t]{0.51\linewidth}
{\large\bfseries\color{model} What I would actually like to ask you}\\[3pt]
\begin{itemize}
  \item Which of those three, if any, is something MURR does routinely --- and which would you
        send elsewhere?
  \item Is there a coupon-exposure path at all: a defined field, a known water chemistry, a
        surface that can be recovered and sectioned?
  \item Is biofilm on cladding-candidate alloys a question anyone in your area is already
        asking, or is it treated as a water-chemistry problem upstream of the materials work?
  \item Would a computed dose field inside a surface film be useful to that work, or is the
        quantity that matters entirely the corrosion outcome?
  \item And the blunt one: is there a version of this that is worth a materials scientist's time,
        or is the honest answer that the interesting part is all microbiology?
\end{itemize}

\vspace{3mm}
\rule{\linewidth}{0.4pt}\\[3pt]
{\small\textbf{One thing I will not do.} There is an organism in the paper, \textit{Ochrobactrum
intermedium} AM7, reclassified in 2020 as \textit{Brucella intermedia}. It is named under both
throughout. No culture has been grown, no strain has been procured, and no containment
determination is asserted anywhere in this work --- containment follows strains, not species, and
it is not mine to declare. If any of this ever moves toward a bench, that determination is the
first step and it belongs to an institution, in a document, with an identifier.}
\end{minipage}

\end{document}
